% ==============================================================================
% FFGFT: Document 247 (English)
% Title: Information, State and Process — FFGFT's Position
% Author: Johann Pascher
% Date: May 2026
% ==============================================================================

\section*{Abstract}

Is information energy? The question is often answered affirmatively in
information physics. FFGFT differentiates fundamentally:

\begin{itemize}
  \item The mere \emph{storage} of information as a state is energy-neutral.
  \item Only \emph{processing} — in particular \emph{erasure} — carries an
        energy dimension.
\end{itemize}

This distinction is not philosophical but empirically decidable: permanent
storage media, magnetic order, and the Casimir effect confirm it directly.
In FFGFT, states are geometric configurations on $T^4$; energy is a derived
quantity that follows from processes, not from the bare existence of structure.

% ------------------------------------------------------------------------------
\section{The Question}
% ------------------------------------------------------------------------------

In information physics — represented by J.~A.~Wheeler's "It from Bit"
\cite{Wheeler1990} and in the IPI context by Denis \cite{Denis2023} —
information is often treated as the ontological primitive from which
energy and matter follow. Landauer's principle \cite{Landauer1961} is
frequently cited as evidence for the physical reality of information:
since erasing a bit costs energy, information must be physical.

This argument conflates two structurally distinct concepts:

\begin{keyresult}[Central Distinction]
\textbf{Information as state} — a configuration, a stored bit, a winding
number — exists without energy expenditure. It is a topological or geometric
property.\\
\textbf{Information as process} — reading, writing, erasing, transforming
— costs energy. Landauer's principle applies \emph{only} to erasure, i.e.\
to a \emph{process}, not to the state itself.
\end{keyresult}

Anyone who conflates these must accept a direct empirical consequence: every
stored bit on a switched-off device would continuously consume energy —
which would make permanent storage media fundamentally impossible. The
existence of such media refutes the equation of information and energy.

% ------------------------------------------------------------------------------
\section{Landauer's Principle — Precisely}
% ------------------------------------------------------------------------------

Rolf Landauer showed in 1961 that the irreversible erasure of one bit
dissipates heat into the surroundings \cite{Landauer1961}. The minimum
energy required is:

\begin{equation}
  E_{\min} = k_B T \ln 2
\end{equation}

where $k_B$ is Boltzmann's constant and $T$ is the absolute temperature.
At room temperature this corresponds to approximately
$2.9 \times 10^{-21}\,\text{J}$ per erased bit.

\textbf{What Landauer shows:} Erasure — an irreversible process — dissipates
energy. This is a consequence of the second law of thermodynamics
\cite{Bennett1982}.

\textbf{What Landauer does not show:} That the bare \emph{existence} of
a stored bit requires energy. A magnetic bit on a hard drive, a winding
number on $T^4$, a spin alignment in a permanent magnet — all exist
without continuous energy input. Only the transition between states costs
energy.

In natural units ($\hbar = c = k_B = 1$), the erasure energy has the
dimension of a temperature, hence a mass, hence an energy — connecting to
$\tilde{T} \cdot m = 1$ in FFGFT. But this dimensional relationship holds
for the \emph{process}, not for the resting state.

% ------------------------------------------------------------------------------
\section{Writing versus Erasing: No Symmetric Energy Requirement}

One might object: writing a bit also costs energy — and if energy is
introduced, mass arises via $E = mc^2$. Would information then carry mass
after all?

This conclusion is wrong for two reasons.

\subsection{Logical Reversibility Decides, Not the Process as Such}

The four possible deterministic operations on one bit are: HOLD, RESET-to-0,
RESET-to-1, and NOT (flip). RESET-to-0 and RESET-to-1 are logically
\emph{irreversible} — the previous state is lost. NOT and HOLD are logically
\emph{reversible} — the previous state can be reconstructed from the result.

Landauer's bound applies \emph{only} to logically irreversible operations,
i.e.\ only to RESET \cite{Landauer1961}. Dago and Bellon demonstrated
experimentally that the NOT operation — writing a new state by bit-flip —
can be carried out in a physically reversible procedure at vanishingly small
energy cost, far below the Landauer bound \cite{Dago2023}. There is no
fundamental energy bound for writing as such.

\subsection{Energy Flows into the Environment, Not into the State}

In practice, writing always costs energy — but this energy does not go
\emph{into the stored bit}. It goes into:

\begin{itemize}
  \item Ohmic heat in the write current (Flash, DRAM)
  \item Magnetic remagnetisation work against the coercive force
  \item Mechanical dissipation in the write head
\end{itemize}

The stored bit itself carries no additional energy or mass. A written
Flash memory does not weigh more than an empty one — even though writing
cost energy. The energy has flowed into the environment, not been stored
in the state.

\subsection{Why Information Does Not Create Mass}

The inference "writing costs energy → $E=mc^2$ → mass arises" commits a
category error. $E=mc^2$ applies to the rest energy of a physical system —
i.e.\ to the mass of the storage medium itself, not to the energy that
dissipates into the environment during the write process.

In FFGFT this is consistent: winding numbers on $T^4$ are topological
invariants. The transition from one configuration to another is a process
with an energy dimension via $\tilde{T}\cdot m=1$ — but the configuration
itself carries no additional rest mass. The process does not create new
mass; it changes the state of a system that already has mass.

\section{Three Empirical Pillars for the Energy-Neutrality of States}

\subsection{Pillar 1: Magnetism}

A permanent magnet is the most tangible evidence for the energy-neutrality
of information states. Spins in a ferromagnetic material are aligned —
an information state in the physical sense (north vs.\ south pole, bit =
direction of magnetisation).

\begin{itemize}
  \item \textbf{Storage:} requires no ongoing energy. The magnet holds
        its alignment without power supply — in principle indefinitely.
  \item \textbf{Remagnetisation:} costs energy. Overcoming the coercive
        field strength requires an external field that performs work.
  \item \textbf{Slow decay:} Over very long timescales, thermal
        fluctuations, cosmic radiation, or other external disturbances
        can alter the alignment \cite{Nattermann1990}. But this decay
        is externally triggered — a process with an external cause,
        not an intrinsic property of the state itself.
\end{itemize}

This is the structural point: even the slow decay of a permanent storage
medium is an externally triggered process, not evidence that the state
itself costs energy.

\subsection{Pillar 2: The Casimir Effect}

The Casimir effect shows that vacuum geometry has physical consequences
without any classical energy store being present. Two uncharged, parallel
metal plates in vacuum attract one another — a measurable force first
predicted in 1948 by Hendrik Casimir \cite{Casimir1948} and first
precisely measured in 1996 by Steve Lamoreaux \cite{Lamoreaux1997}.

The force per unit area is:

\begin{equation}
  P(z) = -\frac{\pi^2}{240}\frac{\hbar c}{z^4}
\end{equation}

where $z$ is the plate separation.

\textbf{What the Casimir effect shows:} The geometric restriction of
the allowed vacuum modes between the plates produces an energy density
difference that generates a force. The \emph{geometry of the vacuum}
— the state — has physical consequences.

\textbf{What the Casimir effect does not show:} That this geometry stores
or consumes energy while the plates are stationary. Only the
\emph{motion of the plates} — the process — performs or absorbs work.

The structural parallel to FFGFT is this: the winding structure on $T^4$ is a
geometric configuration — it exists, has consequences (mass ratios,
coupling constants), but it does not \emph{cost} ongoing energy. Only
dynamic transitions between configurations — processes — carry an energy
dimension. This parallel is an analogy of structure, not an identity of
mechanism.

\subsection{Pillar 3: Permanent Storage Media}

The entire digital civilisation rests on the energy-neutrality of stored
information. Flash memory, magnetic hard drives, optical data carriers —
all use states that persist without power supply.

\begin{itemize}
  \item \textbf{Flash memory} stores bits as charge states in
        floating-gate transistors. The charge is retained without energy
        input — typically for ten years or more.
  \item \textbf{Magnetic tapes} hold data for decades without power.
  \item \textbf{Slow decay:} Quantum tunnelling, thermal activation,
        and ionising radiation can flip bits over very long timescales
        \cite{Bez2003}. But here too: every individual flip is an
        externally triggered process, not a spontaneous property of
        the resting state.
\end{itemize}

If information as a state were not energy-neutral, no permanent storage
medium would be possible. The existence of this technology is the empirical
proof of the central distinction.

% ------------------------------------------------------------------------------
\section{In Natural Units: State, Process, and the Foundational Relation}
% ------------------------------------------------------------------------------

In natural units ($\hbar = c = k_B = 1$), dimensions collapse:

\begin{equation}
  [E] = [m] = [T^{-1}] = [\text{temperature}]
\end{equation}

The erasure energy $E_{\min} = k_B T \ln 2$ becomes in natural units
$E_{\min} = T \ln 2$ — a temperature, a mass, an energy.

FFGFT's foundational relation $\tilde{T} \cdot m = 1$ connects time and
mass directly. A state transition — a process — has a duration and changes
the effective mass of the dynamical field $\Delta m(x,t)$. This is the
natural embedding of Landauer's principle in FFGFT:

\begin{itemize}
  \item \textbf{State:} Winding configuration on $T^4$ — $m$ is
        geometrically fixed by $\xi$, no energy consumption.
  \item \textbf{Process:} Transition between configurations — $\Delta m$
        changes; via $\tilde{T} \cdot m = 1$ this has a time dimension
        and hence an energy dimension.
\end{itemize}

The energy cost of information processing is therefore not an external
addition in FFGFT — it follows from the foundational relation.

% ------------------------------------------------------------------------------
\section{FFGFT's Position}
% ------------------------------------------------------------------------------

\begin{keyresult}[FFGFT and Information]
\textbf{Primitive:} Geometry — $T^4$ winding structure, $\xi$,
$\tilde{T}\cdot m=1$.\\
\textbf{State:} Winding configuration — no energy primitive, analogous
to the magnetic bit or Casimir geometry.\\
\textbf{Process:} State transition — carries an energy dimension,
follows from $\tilde{T}\cdot m=1$ and $\xi$.\\
\textbf{Information:} A descriptive layer over geometric reality, not
an ontological primitive.
\end{keyresult}

Shannon information \cite{Shannon1948} — bits, entropy
$H = -\sum p_j \log p_j$ — is a counting quantity, dimensionless. It
describes states in a probability space. It is not a primitive in FFGFT
because FFGFT does not use probabilities as a foundational concept —
winding numbers are deterministic topological invariants, not random
variables.

% ------------------------------------------------------------------------------
\section{Demarcation}
% ------------------------------------------------------------------------------

\subsection{Wheeler — ``It from Bit''}

John Archibald Wheeler's position \cite{Wheeler1990}: information is
the ontological primitive from which physical reality arises. This places
information before geometry.

FFGFT inverts this: geometry ($T^4$, $\xi$) is the primitive. Information
in the Shannon sense is a description of configurations in this geometry,
not a cause of it.

\subsection{Denis — Landauer Vacuum (IPI)}

Denis \cite{Denis2023} proposes deriving the cosmological constant via
Landauer's principle from an information-theoretic vacuum sector. This
presupposes that information is ontologically fundamental.

FFGFT derives $\Lambda = 1.34 \times 10^{-122}$ (Planck units) directly
from $\xi$ \cite{Pascher_182}, without information theory as an
intermediate layer. The path is shorter and requires no Landauer bridge.

\subsection{RSG — Survival Weights}

RSG \cite{Austin2026} uses survival weights $S_i = e^{-A_i}$ for
histories. This is a probability structure over configurations — a
descriptive layer, not a geometric primitive. From FFGFT's perspective,
RSG is a possible description of selection processes, but not an
explanation of the scale values that determine $\Gamma_i$.

% ------------------------------------------------------------------------------
\section{Symbol List}
% ------------------------------------------------------------------------------

{\small
\adjustbox{max width=\textwidth}{%
\begin{tabular}{p{3.5cm}p{9.5cm}}
\toprule
\textbf{Symbol} & \textbf{Meaning} \\
\midrule
$k_B$ & Boltzmann constant: $k_B = 1.380649 \times 10^{-23}\,\text{J/K}$ \\
$T$ & Absolute temperature (Kelvin) \\
$E_{\min} = k_B T \ln 2$ & Landauer bound: minimum energy to erase one bit \\
$H = -\sum p_j \log p_j$ & Shannon entropy; dimensionless in bits (base 2) \cite{Shannon1948} \\
$\xi = 4/30000$ & FFGFT coupling parameter; dimensionless \\
$\tilde{T} \cdot m = 1$ & FFGFT foundational relation (natural units) \\
$T^4$ & 4-dimensional torus; FFGFT ontological primitive \\
$\Delta m(x,t)$ & Dynamical mass field in T0; carries state transitions \\
$P(z) = -\frac{\pi^2}{240}\frac{\hbar c}{z^4}$ & Casimir pressure; $z$ = plate separation \cite{Casimir1948} \\
$\Lambda = 1.34 \times 10^{-122}$ & Cosmological constant from $\xi$ (Planck units) \cite{Pascher_182} \\
\bottomrule
\end{tabular}%
}
}

% ------------------------------------------------------------------------------
\section*{References}
% ------------------------------------------------------------------------------

\begin{thebibliography}{99}

\bibitem{Landauer1961}
R.~Landauer,
\textit{Irreversibility and Heat Generation in the Computing Process},
IBM Journal of Research and Development \textbf{5}(3), 183--191 (1961).
DOI: 10.1147/rd.53.0183.

\bibitem{Bennett1982}
C.~H.~Bennett,
\textit{The Thermodynamics of Computation — A Review},
International Journal of Theoretical Physics \textbf{21}(12), 905--940 (1982).
DOI: 10.1007/BF02084158.

\bibitem{Shannon1948}
C.~E.~Shannon,
\textit{A Mathematical Theory of Communication},
Bell System Technical Journal \textbf{27}(3), 379--423 (1948).
DOI: 10.1002/j.1538-7305.1948.tb01338.x.

\bibitem{Casimir1948}
H.~B.~G.~Casimir,
\textit{On the Attraction Between Two Perfectly Conducting Plates},
Proceedings of the Koninklijke Nederlandse Akademie van Wetenschappen
\textbf{51}, 793--795 (1948).

\bibitem{Lamoreaux1997}
S.~K.~Lamoreaux,
\textit{Demonstration of the Casimir Force in the 0.6 to 6 $\mu$m Range},
Physical Review Letters \textbf{78}(1), 5--8 (1997).
DOI: 10.1103/PhysRevLett.78.5.

\bibitem{Wheeler1990}
J.~A.~Wheeler,
\textit{Information, Physics, Quantum: The Search for Links},
in: W.~H.~Zurek (ed.), \textit{Complexity, Entropy, and the Physics of
Information}, Addison-Wesley, Redwood City (1990), pp.~3--28.

\bibitem{Nattermann1990}
T.~Nattermann,
\textit{Theory of the Random Field Ising Model},
in: A.~P.~Young (ed.), \textit{Spin Glasses and Random Fields},
World Scientific, Singapore (1998), pp.~277--298.

\bibitem{Bez2003}
R.~Bez et al.,
\textit{Introduction to Flash Memory},
Proceedings of the IEEE \textbf{91}(4), 489--502 (2003).
DOI: 10.1109/JPROC.2003.811702.

\bibitem{Denis2023}
O.~Denis,
\textit{Informational Nature of Dark Matter and Dark Energy and the
Cosmological Constant},
IPI Letters \textbf{1}, 66--77 (2023).
DOI: 10.59973/ipil.36.

\bibitem{Austin2026}
P.~M.~Austin,
\textit{Recursive Survival Geometry: A Conditional Recursive Formalism
and Analogue-Media Test Protocol for Survival-Weighted Representation},
v1.06, Information Physics Institute (2026).

\bibitem{Dago2023}
S.~Dago and L.~Bellon,
\textit{Logical and Thermodynamical Reversibility: Optimized Experimental
Implementation of the NOT Operation},
arXiv:2302.11908 (2023).

\bibitem{Pascher_182}
J.~Pascher,
\textit{T0-Universe Maximum Scale — Cosmological Constant and Hubble
Parameter from $\xi$},
Document 182, FFGFT Document Series (2026).
GitHub: jpascher/T0-Time-Mass-Duality.
Zenodo: DOI 10.5281/zenodo.20474821.

\end{thebibliography}
